\section{Considerações Finais}
\label{sec:consideracoes-finais}
\label{sec:con}

Este trabalho apresentou o desenvolvimento do \textit{DyraSQL}, um mecanismo de decisão para roteamento dinâmico de consultas SQL baseado em metadados de tabelas \textit{Apache Iceberg}. O problema abordado consistiu na necessidade de arquiteturas que se adaptem ao tipo de consulta, evitando provisionamento excessivo de infraestrutura robusta para consultas leves e garantindo recursos adequados para consultas pesadas, resultando em melhor custo-benefício no processamento.

Os objetivos propostos foram alcançados através da análise de metadados \textit{Apache Iceberg}, estabelecimento de regras de decisão para roteamento, desenvolvimento de um \textit{framework} funcional em ambiente local e avaliação da lógica de roteamento. A metodologia envolveu análise dos metadados disponíveis, definição de critérios baseados em volume de dados, complexidade da consulta e histórico de execuções, e verificação da lógica proposta através de testes com tabelas \textit{Iceberg} geradas localmente.

Os principais resultados demonstraram que a lógica de roteamento funciona de forma consistente, sendo que o \textit{framework} é capaz de calcular \textit{scores} e direcionar consultas para os \textit{clusters} apropriados baseado nas características dos dados e na complexidade das consultas. A análise dos metadados revelou informações sobre volume de dados e complexidade que podem ser utilizadas para determinar a arquitetura mais adequada, sendo que o uso do \textit{EXPLAIN (TYPE IO)} do \textit{Trino} fornece estimativas suficientes para o Algoritmo de decisão. A avaliação confirmou que os pesos do Algoritmo funcionam de forma adequada, sendo que a combinação dos fatores de volume, complexidade e histórico resulta em \textit{scores} que diferenciam consultas leves, intermediárias e pesadas.

As contribuições deste trabalho incluem um \textit{framework} para análise de metadados \textit{Apache Iceberg} para determinação de arquiteturas dinâmicas baseado em volume de dados, complexidade de consultas e histórico de execuções; regras de decisão para roteamento automático através de algoritmo de pontuação que combina múltiplos fatores; um protótipo funcional integrando \textit{Trino}, catálogo \textit{Iceberg} REST, MinIO e PostgreSQL, que direciona consultas para \textit{clusters small}, \textit{medium} ou \textit{large}; e a avaliação da lógica de roteamento através de testes com dados sintéticos em ambiente local.

Como limitações do estudo, destacam-se o uso de \textit{dataset} sintético em escala inferior à de produção e a dependência de histórico suficiente no PostgreSQL, sendo que consultas sem histórico apresentam maior probabilidade de decisões menos precisas. A extração de metadados apresenta limitações para consultas muito complexas, sendo que o tempo do \textit{EXPLAIN (TYPE IO)} pode impactar a latência total. O sistema também depende da qualidade dos metadados do \textit{Apache Iceberg}, sendo que tabelas com metadados desatualizados podem resultar em estimativas imprecisas do volume efetivo.

\subsection{Trabalhos Futuros}

Trabalhos futuros concentram-se na aplicação do \textit{DyraSQL} em ambientes de nuvem. O núcleo do sistema, composto pelo cálculo do \textit{Score}, pelo \textit{Trino Gateway} e pelas três faixas de \textit{cluster}, permanece o mesmo, sendo que o que se altera é a implantação. \textit{Clusters} gerenciados em nuvem, por exemplo EMR \cite{awsemr2025}, e o armazenamento de objetos do provedor substituiriam o MinIO e as instâncias \textit{Trino} locais, permitindo avaliar o roteamento com volumes de produção e com elasticidade típica de plataformas gerenciadas.

Outra direção consiste em substituir os pesos $w_1$, $w_2$ e $w_3$, atualmente fixos, por um modelo de aprendizado supervisionado treinado sobre o histórico acumulado no PostgreSQL, de forma similar ao Amazon Redshift MDDL \cite{ding2024}, além de avaliar o \textit{DyraSQL} sobre tabelas de centenas de gigabytes a terabytes, verificando se as faixas de \textit{score} permanecem adequadas em escala real.

Uma limitação da arquitetura atual é que os três \textit{clusters} são \textit{containers} persistentes com capacidade fixa, de modo que um pico de consultas concorrentes na mesma faixa de \textit{score} pode esgotar o \textit{cluster} de destino mesmo com os demais subutilizados. Um trabalho futuro é migrar para um modelo orientado a \textit{tasks}, análogo à distinção que o Amazon ECS estabelece entre \textit{services} de longa duração e \textit{tasks} avulsas sob demanda \cite{awsecs2025}, dimensionando o ambiente de execução individualmente por consulta em vez de compartilhar um \textit{cluster} fixo. Também é necessário definir uma estratégia de \textit{fallback} para os casos em que a análise de metadados não é concluída (indisponibilidade do \textit{Trino} ou \textit{timeout}), direcionando a um \textit{cluster} padrão em vez de retornar erro ao cliente — situação não observada nos testes, mas esperada antes de uma implantação em produção.
